\section{Discussion}

\subsection{Operational Implications}

The temporal results change the operational case for hygiene-augmented scoring. HygienePrio~\cite{paper4hygieneprio} showed that HygienePrio outperforms EPSS-only at a single window. This paper shows that EPSS-only decays sharply across rolling maintenance windows under the simulated bi-weekly cadence: by Window~3 EPSS-only's mean P@50 ($0.087$) is comparable to the Random baseline ($0.084$), and by Window~6 EPSS-only ($0.034$) is essentially indistinguishable from random ($0.054$). HRS-only, the simplest hygiene baseline, overtakes EPSS-only at Window~2 and dominates it for every subsequent window.

The concrete operational recommendation is that hygiene signals are not supplementary to EPSS --- they are necessary for sustained prioritization quality. EPSS provides strong initial triage (Window~1 gap $= 26.4$ pp between HygienePrio-full and EPSS-only), but its advantage over hygiene-based baselines is restricted to Window~1 alone. The HygienePrio-vs-EPSS gap widens to $46.5$ pp at Window~6 because HygienePrio-full holds P@50 near $0.5$ while EPSS-only collapses toward the random floor.

\subsection{Connection to the Research Sequence}

Prior single-window evaluations~\cite{paper1vulnprio,paper3hygienebench,paper4hygieneprio} were all single-window studies. This paper reveals that the choice of evaluation window is a major moderating variable that prior studies did not account for. VulnPrio's null result (CVE-level augmentation $\approx$ EPSS-only) was measured at a single window when EPSS still has maximal signal. Had it been measured at Window~3+, EPSS-only would have been essentially random and any host-aware augmentation would have appeared superior. This temporal lens retroactively reframes VulnPrio: its null result describes the regime in which EPSS is strong, not the regime in which EPSS-only would be operated by an actual remediation team across months.

\subsection{Why HRS Does Not Collapse}

The asymmetric collapse pattern is not an artifact of the ground-truth definition. HRS does not collapse because hygiene posture is a property of the host and changes only when remediation acts on it. A host with high patch debt remains a high-debt host across multiple windows unless many of its pairs are selected; under capacity constraints, only a fraction of high-HRS hosts are touched per window. By contrast, EPSS is a property of the CVE, and remediating a single high-EPSS pair removes that CVE from the unpatched backlog for the entire fleet. The compositional asymmetry between host-level state (slow to drain) and CVE-level state (drained by any acting host) is the structural reason hygiene signals are temporally resilient.

\subsection{Limitations}

\textbf{Simplified remediation model.} Patch application reduces patch debt by a fixed 15\% per window, which may not reflect real-world patch-velocity heterogeneity.

\textbf{Fixed remediation policy.} All methods are evaluated on the fleet trajectory generated by HygienePrio-full's selections. This avoids non-identifiability but means the comparison is not a true counterfactual --- an EPSS-only-driven fleet would evolve differently, and the absolute P@50 numbers in later windows would shift.

\textbf{Window-recomputed ground truth.} The HRS-percentile threshold is recomputed each window, which interacts with the patch-application dynamics. A fixed Window~1 threshold was considered but would have produced trivially decreasing positive counts; the per-window threshold preserves the statistical interpretability of P@50 across windows at the cost of a structural advantage for HRS-based methods that is acknowledged in \S\ref{sec:threats}.

\textbf{Synthetic only.} All results are bounded to the synthetic EEHDA evaluation context. External validation on real fleet data with longitudinal exploitation outcome data is required before any deployment recommendation.
